\documentclass[11pt]{article}

% --- packages (all standard in a basic TeX Live install) ---
\usepackage[utf8]{inputenc}
\usepackage[margin=1in]{geometry}
\usepackage{amsmath,amssymb}
\usepackage{booktabs}
\usepackage{xcolor}
\usepackage{listings}
\usepackage{enumitem}
\usepackage[colorlinks=true,linkcolor=blue,urlcolor=blue]{hyperref}

\lstset{
  basicstyle=\ttfamily\small,
  breaklines=true,
  frame=single,
  columns=fullflexible,
  backgroundcolor=\color{black!3},
  keepspaces=true,
}

\newcommand{\code}[1]{\texttt{#1}}
\newcommand{\R}{\mathbb{R}}

\title{\textbf{Cross-Modal Brain MRI Retrieval}\\[2pt]
\large A Classical, Registration-Based Pipeline (public score $0.904$)}
\author{EHL Paris --- Medical Image Retrieval}
\date{2026-06-27}

\begin{document}
\maketitle

\begin{abstract}
We describe, from scratch, the retrieval system that reaches a public score of
$0.90444$ on the \emph{EHL Paris: Medical Image Retrieval} challenge, up from a
$0.651$ baseline. The task is to match a T1 post-contrast query brain MRI to its
same-subject T2 target across three datasets with increasing geometric difficulty.
The final system is fully classical (no deep learning, no augmentation): per-volume
intensity normalization, an \emph{intramodal template normalization} step that
removes pose, a PCA + Ridge cross-modal embedding learned on the only labelled
dataset, cosine scoring, and a one-to-one Hungarian assignment that exploits the
bijective structure of each retrieval pool. A key methodological contribution is a
\emph{synthetic-dataset2 validation harness} that lets us rank candidate methods
offline, without spending the limited Kaggle submission budget.
\end{abstract}

\tableofcontents

% ======================================================================
\section{Problem statement}

For each \emph{query} volume (a 3D T1 post-contrast MRI) we must rank all
\emph{target} volumes (3D T2 MRI) of a gallery so that the true same-subject target
is ranked as high as possible. All images are NIfTI volumes in RAS orientation at
$1\,\text{mm}$ isotropic spacing; shapes vary (e.g. $240\times240\times155$).

The data is split into three independent pools:
\begin{itemize}[nosep]
  \item \textbf{dataset1}: registered preoperative pairs. The \emph{only} labelled
        training set ($350$ pairs). Query and target of a pair share a common grid.
  \item \textbf{dataset2}: same source as dataset1, but each query and target is
        \emph{independently} warped by a random rigid transform (rotation +
        translation) \emph{and} a nonlinear (elastic) deformation. No labels.
  \item \textbf{dataset3}: preoperative$\to$intraoperative pairs. The intraoperative
        target is resampled into the query's physical space, so the pair already
        shares geometry, but surgery shifts/removes tissue. No labels.
\end{itemize}

\paragraph{Counts.} dataset1: $40/100$ val/test queries (and equal galleries);
dataset2: $40/100$; dataset3: $20/77$. A full submission has $377$ rows.

% ======================================================================
\section{Evaluation metric and the bijection insight}

The score is the mean reciprocal rank (MRR), computed per dataset and averaged:
\begin{equation}
\text{score} = \tfrac{1}{3}\big(\text{MRR}_{1}+\text{MRR}_{2}+\text{MRR}_{3}\big),
\qquad
\text{MRR}_d = \frac{1}{|Q_d|}\sum_{q\in Q_d}\frac{1}{\text{rank}_q},
\end{equation}
where $\text{rank}_q$ is the position of the true target in the ranking for query
$q$. A partial submission of one dataset reports $\text{MRR}_d/3$, so multiplying by
$3$ recovers that dataset's MRR --- the basis of all our diagnostics.

\paragraph{Bijection.} In every pool the number of queries equals the number of
gallery targets, and the matching is one-to-one. Retrieval is therefore exactly a
\emph{linear assignment problem}: find the permutation $\pi$ matching queries to
targets. If a similarity matrix $S$ ranks the right target first for all queries,
$\text{MRR}=1$. This motivates the Hungarian re-ranking in
Section~\ref{sec:hungarian}.

% ======================================================================
\section{Preprocessing: intensity normalization}\label{sec:prep}

Each volume $V$ is loaded and reduced to a normalized intensity field on a fixed
$G\times G\times G$ grid (we use $G=44$).

\begin{enumerate}[nosep]
  \item Load NIfTI; if 4D keep the first channel; replace NaN/Inf with $0$.
  \item Foreground mask $M=\{|V|>10^{-6}\}$.
  \item Robust window: $lo,hi$ = the $1$st and $99$th percentiles of $V$ over $M$.
  \item Clip-and-scale: $V \leftarrow \mathrm{clip}\!\big((V-lo)/(hi-lo),\,0,\,1\big)$,
        and set background ($\lnot M$) to $0$.
  \item Resample $V$ onto the $G^3$ grid by trilinear interpolation
        (\code{scipy.ndimage.map\_coordinates}, order $1$).
\end{enumerate}

This makes intensities comparable across scans of differing dynamic range, and the
fixed grid makes every volume a vector of the same length $G^3$.

% ======================================================================
\section{Cross-modal embedding: PCA + Ridge}\label{sec:embed}

Query (T1) and target (T2) are different modalities, so raw intensities are not
directly comparable. We learn the cross-modal relationship on the $350$ labelled
dataset1 pairs.

Let $\phi(\cdot)\in\R^{G^3}$ be the flattened, mean-centred, $L_2$-normalized grid
feature of a (geometry-normalized) volume. From the training pairs
$\{(\phi^q_i,\phi^t_i)\}_{i=1}^{350}$:

\begin{enumerate}[nosep]
  \item Fit a whitening PCA on the queries, $P_q:\R^{G^3}\!\to\R^{c}$, and a separate
        whitening PCA on the targets, $P_t$, with $c=\min(128,\,n-1)$.
  \item Fit a ridge regression $W$ mapping query codes to target codes:
        \begin{equation}
          W=\arg\min_{W}\ \sum_i \big\|W\,P_q(\phi^q_i)-P_t(\phi^t_i)\big\|^2
          +\alpha\|W\|_F^2,\qquad \alpha=100 .
        \end{equation}
\end{enumerate}

At inference, a query embeds as $u = \mathrm{norm}\!\big(W\,P_q(\phi^q)\big)$ and a
target as $v = \mathrm{norm}\!\big(P_t(\phi^t)\big)$; their similarity is the cosine
$u^\top v$. Intuitively, PCA denoises and decorrelates each modality, and ridge
learns the linear T1$\to$T2 code transformation, so matched subjects land nearby in
the shared target code space.

% ======================================================================
\section{Geometry: template normalization}\label{sec:template}

The embedding of Section~\ref{sec:embed} assumes query and target occupy a
\emph{common geometric frame}. That holds for dataset1 (registered) but not for
dataset2, where pose differs per volume. We restore the common frame by
\emph{template normalization}.

\paragraph{Templates.} Because dataset1 pairs are registered, the voxel-wise means
\begin{equation}
T_1=\frac{1}{350}\sum_i V^q_i,\qquad
T_2=\frac{1}{350}\sum_i V^t_i
\end{equation}
are two sharp templates that, crucially, \emph{share one grid}: $T_1$ (a mean T1
brain) and $T_2$ (a mean T2 brain) are mutually registered.

\paragraph{Registration.} Each query is rigidly aligned to $T_1$ and each target to
$T_2$ --- an \emph{intramodal} registration (same modality on both sides), which is
far more robust than cross-modal alignment. We optimize a 6-parameter rigid
transform $\theta=(\text{rot}_x,\text{rot}_y,\text{rot}_z,t_x,t_y,t_z)$ to maximize
the normalized cross-correlation (NCC) over the foreground:
\begin{equation}
\theta^\star=\arg\max_{\theta}\ \mathrm{NCC}\big(T_m,\ R_\theta[V]\big),
\end{equation}
where $R_\theta$ applies the rigid transform about the volume centre. Optimization
uses Powell's method from $4$ starts (identity $+$ three seeded rotations) on a
coarse $20^3$ copy for speed; the recovered $\theta^\star$ is then applied at full
grid resolution. The aligned volume feeds $\phi(\cdot)$ as before.

\paragraph{Why it works.} After normalization, query and target of the \emph{same}
subject occupy the same template frame, so the dataset1-trained map of
Section~\ref{sec:embed} transfers to dataset2. The cost is $O(N)$ (one registration
per volume), not $O(N^2)$ (every query-target pair). This single step raised
dataset2 MRR from $\sim0.39$ to $\sim0.75$ and dataset1 from $\sim0.877$ to
$\sim0.964$.

\paragraph{When NOT to register (dataset3).} dataset3 targets are already resampled
into the query's space, so query and target \emph{already share geometry}.
Registering an intraoperative volume to a \emph{preoperative} template only injects
misalignment around surgical changes. Empirically, template normalization
\emph{hurts} dataset3 (MRR $0.690\to0.251$). Instead we skip registration and feed
the raw grid feature $\phi(V)$ directly. The grid feature itself is far stronger
than the legacy hand-crafted features, lifting dataset3 to $\text{MRR}\approx1.0$.

% ======================================================================
\section{One-to-one Hungarian assignment}\label{sec:hungarian}

Given the cosine similarity matrix $S\in\R^{N\times N}$ for a pool, we exploit the
bijection (Section~2). We solve the optimal assignment
\begin{equation}
\pi^\star=\arg\max_{\pi\in \mathfrak{S}_N}\ \sum_{i} S_{i,\pi(i)}
\end{equation}
with the Hungarian algorithm (\code{scipy.optimize.linear\_sum\_assignment} on
$-S$). For each query we then place its assigned target first (by boosting
$S_{i,\pi^\star(i)}$ above all other entries) and rank the remainder by raw
similarity. This converts a globally consistent matching into per-query rankings and
gave a large, repeatable gain (e.g.\ $0.557\to0.594$ early on).

% ======================================================================
\section{The synthetic-dataset2 validation harness}\label{sec:harness}

Kaggle allows only $100$ submissions/day and dataset2/3 have no labels, so blind
iteration is expensive. We build a \emph{local} proxy for dataset2 from the labelled
dataset1.

\paragraph{Construction.} Split the $350$ dataset1 pairs into a clean \emph{train}
set (used to fit templates and the PCA/Ridge map) and an \emph{eval} set. Each eval
query and target is warped \emph{independently} by:
\begin{itemize}[nosep]
  \item a random rigid transform (rotation up to $\pm r$ degrees, translation up to
        $\pm s$ voxels), and
  \item an elastic deformation: a Gaussian-smoothed ($\sigma$) random displacement
        field scaled by $\alpha$, applied via \code{map\_coordinates}.
\end{itemize}
A candidate method produces a similarity matrix on the eval set, and we report
Hungarian-recovery MRR (the eval diagonal is ground truth).

\paragraph{Calibration.} We tuned $(r,s,\sigma,\alpha)=(12,6,8,3)$ so that the
existing canonical baseline scores $\approx0.26$ on the harness, matching its real
dataset2 MRR. At that operating point template normalization scored $\approx0.55$
--- correctly predicting the real win \emph{before} any submission.

\paragraph{Scope and a caveat.} The harness is reliable for \emph{relative ranking}
of methods under dataset2-style distortion. It is \emph{not} reliable for absolute
choices on real dataset1/3, nor for resolution choices: it favored grid $56$
($0.72$ vs $0.56$), but on real data grid $56$ \emph{regressed} both d1 and d2
(likely PCA overfitting with only $350$ training pairs). Grid $44$ is optimal.

% ======================================================================
\section{Full pipeline, per dataset}

\begin{center}
\begin{tabular}{@{}llll@{}}
\toprule
Dataset & Geometry handling & Feature & Map + rank \\
\midrule
dataset1 & template normalization ($T_1/T_2$) & $\phi$ on $44^3$ & PCA+Ridge, cosine, Hungarian\\
dataset2 & template normalization ($T_1/T_2$) & $\phi$ on $44^3$ & PCA+Ridge, cosine, Hungarian\\
dataset3 & \emph{none} (already aligned)       & $\phi$ on $44^3$ & PCA+Ridge, cosine, Hungarian\\
\bottomrule
\end{tabular}
\end{center}

\noindent End-to-end for a pool:
\begin{enumerate}[nosep]
  \item Build $T_1,T_2$ and fit $P_q,P_t,W$ on dataset1 (once).
  \item For each query/target: normalize intensities $\to$ resample to $44^3$ $\to$
        (register to its template, \emph{except} dataset3) $\to$ $\phi$.
  \item Embed: $u=\mathrm{norm}(W P_q(\phi^q))$, $v=\mathrm{norm}(P_t(\phi^t))$.
  \item $S=UV^\top$; Hungarian boost; write space-separated ranking per query.
\end{enumerate}

% ======================================================================
\section{Results}

\begin{center}
\begin{tabular}{@{}lcc@{}}
\toprule
Submission & Public score & Note \\
\midrule
canon20 flipaug (prior best)                  & 0.65137 & session start \\
d1 pca + d2 \textbf{template} + d3 pca          & 0.77071 & d2 fix \\
d1 \textbf{template} + d2 template + d3 pca     & 0.80127 & d1 also helped by template \\
d1 template + d2 template + d3 \textbf{grid}    & \textbf{0.90444} & \textbf{best; goal hit} \\
grid56 variants                                & 0.850 / 0.866 & regressed (see \S\ref{sec:harness}) \\
\bottomrule
\end{tabular}
\end{center}

\noindent Approximate per-dataset MRR of the best mix:
$\text{MRR}_1\!\approx\!0.964$, $\text{MRR}_2\!\approx\!0.749$,
$\text{MRR}_3\!\approx\!1.0$, giving $(0.964+0.749+1.0)/3\approx0.904$.

% ======================================================================
\section{Negative results}

The harness and partial submissions ruled out several directions:
\begin{itemize}[nosep]
  \item \textbf{Pose-invariant descriptors} (intensity/gradient histograms, radial
        profiles, shape eigenvalues): too lossy; discard the spatial structure that
        distinguishes subjects ($\approx0.11$ on the harness).
  \item \textbf{Pairwise rigid edge-NCC re-rank}: elastic warp defeats rigid
        alignment and the optimizer hits local minima ($0.20<0.26$).
  \item \textbf{grid $56$}: better on the harness, worse on real data.
  \item \textbf{Deep contrastive 3D model on $850$ geom+contrast augmented pairs}
        (ROCm GPU): holdout MRR $\approx0.04$, far below classical $0.90$. The
        scratch model does not learn competitive embeddings even with augmentation.
  \item \textbf{BrainIAC foundation embeddings} (cosine / adapter / patch): all
        below the classical pipeline (historical).
\end{itemize}

% ======================================================================
\section{Reproduction}

\begin{lstlisting}[language=bash]
# Best dataset2 (template normalization) + dataset1
python d2_template_retrieval.py --datasets dataset1 dataset2 \
  --grid 44 --assignment --out submissions/d1d2_template_g44_hungarian.csv

# Best dataset3 (grid feature, NO registration)
python d2_template_retrieval.py --datasets dataset3 \
  --grid 44 --no-register --assignment \
  --out submissions/d3_gridfeat_g44_hungarian.csv

# Offline method validation (no Kaggle needed)
python synthetic_d2_eval.py --n-eval 60 --n-train 250 --grid 44 \
  --max-rot-deg 12 --max-shift 6 --elastic-sigma 8 --elastic-alpha 3 \
  --methods canonical template

# Validate a submission CSV before sending
python classical_retrieval.py validate submissions/<file>.csv
\end{lstlisting}

\paragraph{Key source files.}
\code{d2\_template\_retrieval.py} (production pipeline; caches normalized features
in \code{.d2cache/}); \code{synthetic\_d2\_eval.py} and \code{d2\_methods.py} (the
validation harness and candidate scorers); \code{assignment\_rerank.py} (Hungarian
re-ranking of classical scores); \code{classical\_retrieval.py} (feature cache and
submission validation).

\paragraph{Compile.}
\begin{lstlisting}[language=bash]
pdflatex architecture.tex
\end{lstlisting}

\end{document}
